%%%%%%%%%%%%%%%%%%%%%%%%%%%%%%%%%%%%%%%%%%%%%%%%%%%%%%%%%%%%%%%%%%%%%%%%%%%%%%%%
%                                                                              %
% BlackArch Linux Guide                                                        %
%                                                                              %
%%%%%%%%%%%%%%%%%%%%%%%%%%%%%%%%%%%%%%%%%%%%%%%%%%%%%%%%%%%%%%%%%%%%%%%%%%%%%%%%

\documentclass[a4paper, oneside, 11pt]{book}

%%% INCLUDES %%%
\renewcommand{\familydefault}{\sfdefault}

\usepackage{array}
\usepackage{color}
\usepackage{enumerate}
\usepackage{fancyhdr}
\usepackage{fancyvrb}
\usepackage{geometry}
\usepackage{graphicx}
\usepackage{html}
\usepackage{hyperref}
\usepackage{ifpdf}
\usepackage{listings}
\usepackage{pstricks}
\usepackage{supertabular}
\usepackage{tocloft}
\usepackage[utf8]{inputenc}

%%% LAYOUT %%%
\setlength{\parindent}{0em}
\setlength{\parskip}{1.5ex plus0.5ex minus0.5ex}
\geometry{left=2.5cm, textwidth=16cm, top=3cm, textheight=25cm, bottom=3cm}
\widowpenalty=2000
\clubpenalty=1000
\frenchspacing
\sloppy
\pagecolor[HTML]{FFFFFF}
\color[HTML]{333333}
\setcounter{tocdepth}{10}
\setcounter{secnumdepth}{10}

\hypersetup{
  pdftitle={BlackArch Linux, Guia BlackArch Linux},
  pdfsubject={BlackArch Linux, Guia BlackArch Linux},
  pdfauthor={BlackArch Linux, BlackArch Linux},
  pdfkeywords={BlackArch Linux, Penetration Testing, Security, Hacking, Linux},
  pdfcenterwindow=true,
  colorlinks=true,
  breaklinks=true,
  linkcolor=red,
  menucolor=red,
  urlcolor=red
}

% syntax highlighting
% all options should be set here document wide
\lstset{
backgroundcolor=\color[HTML]{EEEEEE},
frame=single,
basicstyle=\footnotesize\ttfamily,
float,
deletekeywords={return},
otherkeywords={mkdir, curl, sudo, sha1sum, grep, cut, sort, wget, makepkg,
pacman, blackman},
keywordstyle=\color{orange},
commentstyle=\color{blue},
stringstyle=\color{red},
language=bash,
showspaces=false,
showtabs=false,
tabsize=2
}

%%% HEADER / FOOTER %%%
\setlength{\headheight}{33pt}
\setlength{\headsep}{33pt}
\lhead{{\includegraphics[width=1cm,height=1cm]{images/logo.png}}}
\rhead{Guia BlackArch Linux}

%%% CUSTOM MACROS %%%
% for html links
\ifpdf\else
\def\href#1#2{\htmladdnormallink{#2}{#1}}
\fi

%------------------%
%  TITLE PAGE      %
%------------------%
\begin{document}
\pagestyle{empty}
\begin{center}
\begin{figure}[htbp]
\centering
\vspace{0.5cm}
\includegraphics[width=8cm]{images/logo.png}
\label{fig:logo}
\end{figure}
\vspace{0.5cm}
\Huge{\textbf{Guia BlackArch Linux}}\\
\vspace{1cm}
\Large{\color{red}https://www.blackarch.org/}\\
\vspace{0.5cm}
\end{center}
\newpage
\tableofcontents
\newpage
\pagestyle{fancy}

%------------------%
%  Chapter 1       %
%------------------%

\chapter{Introdução}

\section{Resumo}
O guia do BlackArch linux é dividido em algumas partes:
\begin{itemize}
\item Introdução - Apresenta uma visão geral e informações adicionais sobre o projeto
\item Guia do Usuário - Tudo que o usuário precisa saber para usar o BlackArch de forma efetiva
\item Guia do Desenvolvedor - Como começar a desenvolver e contribuir para o BlackArch
\item Guia das Ferramentas - Detalhes e exemplos de uso das ferramentas
\end{itemize}

\section{O que é o BlackArch Linux?}
BlackArch é uma distribuição Linux para testes de penetração e pesquisadores de segurança.É uma derivação do \href{https://www.archlinux.org/}{ArchLinux} e os usuários podem instalar componentes do BlackArch individualmente ou por grupos em cima da distribuição.
As ferramentas são distribuídas através do Arch Linux
\href{https://wiki.archlinux.org/index.php/Unofficial\_User\_Repositories}
{unofficial user repository} então você pode instalar o BlackArch em cima de uma instalação Arch Linux já existente.Os pacotes pode ser instalados individualmente ou por categorias.

O repositório continua aumentando e já possui  \href{https://www.blackarch.org/tools.html}{1300} ferramentas.
Todas as ferramentas são testadas depois da adição de algum codigo mantendo a qualidade do repositório.
% should quickly describe the testing methods/code review procedures etc

\section{A historia do BlackArch Linux}
Em breve...

\section{Plataformas suportadas}
Em breve...

\section{Junte-se ao BlackArch}
Você pode se comunicar com a equipe BlackArch através das formas abaixo:

Website: \url{https://www.blackarch.org/}

Mail: \href{mailto:team@blackarch.org}{team@blackarch.org}

IRC: \url{irc://irc.freenode.net/blackarch}

Twitter: \url{https://twitter.com/blackarchlinux}

Github: \url{https://github.com/Blackarch/}

%------------------%
%  Chapter 2       %
%------------------%


\chapter{Guia do Usuário}

\section{Instalação} 
A seção seguinte mostra como configurar o repositório BlackArch e instalar os pacotes.
O BlackArch suporta a instalação direta pelo repositório usando pacotes binários ou compilando, como também através de outras fontes.
O BlackArch é compatível como toda instalação do Arch.A instalação é feita através de um repositório não oficial, se você quiser usar uma ISO olhe na seção \href{https://www.blackarch.org/downloads.html#iso}{Live ISO}.

\subsection{Instalando em cima do ArchLinux}
Execute \href{https://blackarch.org/strap.sh}{strap.sh} como root e siga as instruções abaixo.
\begin{lstlisting}
   curl -O https://blackarch.org/strap.sh
   sha1sum strap.sh # Deve ser igual a: 86eb4efb68918dbfdd1e22862a48fda20a8145ff
   sudo ./strap.sh
\end{lstlisting}
Agora baixe uma copia atualizada da lista de pacotes e os sincronize.
\begin{lstlisting}
  sudo pacman -Syyu
\end{lstlisting}


\subsection{Instalando os pacotes}
Agora você pode instalar os pacotes do repositório do blackarch.
\begin{enumerate} 
\item Para listar todas as ferramentas disponíveis execute
\begin{lstlisting}
  pacman -Sgg | grep blackarch | cut -d' ' -f2 | sort -u
\end{lstlisting}

\item Para instalar todas a ferramentas execute
\begin{lstlisting}
  pacman -S blackarch
\end{lstlisting}

\item Para instalar toda uma categoria execute
\begin{lstlisting}
  pacman -S blackarch-<category>
\end{lstlisting}

\item Para ver as categorias execute
\begin{lstlisting}
  pacman -Sg | grep blackarch
\end{lstlisting}

\end{enumerate}

\subsection{Instalando os pacotes pelo codigo fonte}
De maneira alternativa você pode criar os pacotes do BlackArch diretamente do codigo fonte.Você pode achar os PKGBUILDs no
\href{https://github.com/BlackArch/blackarch/tree/master/packages}{github}.Para criar todo o repositório você pode usar a ferramenta
\href{https://github.com/BlackArch/blackman}{Blackman}.
\begin{itemize}
\item Primeramente, se você tem que instalar o Blackman, se o repositório do BlackArch está configurado na sua máquina, você pode instalar o Blackman:
\begin{lstlisting}
  pacman -S blackman
\end{lstlisting}

\item Você pode construir e instalar o Blackman pelo código fonte:
\begin{lstlisting}
  mkdir blackman
  cd blackman
  wget https://raw2.github.com/BlackArch/blackarch/master/packages/blackman/PKGBUILD
  # Tenha certeza que o PKGBUILD não possui codigo malicioso.
  makepkg -s
\end{lstlisting}

\item Ou você pode instalar o Blackman pelo AUR::
\begin{lstlisting}
  <quaquer ajudante do AUR > -S blackman
\end{lstlisting} 

\end{itemize}

\subsection{Uso básico do Blackman} Blackman possui um uso simples, suas flags são um pouco diferentes do pacman.O básico será mostrado abaixo.
\begin{itemize}
\item Baixa, compila e instala o pacote:
\begin{lstlisting}
  sudo blackman -i pacote
\end{lstlisting}

\item Baixa, compila e instala toda uma categoria:
\begin{lstlisting}
  sudo blackman -g grupo
\end{lstlisting}

\item Baixa, compila e instala todas as ferramentas do BlackArch:
\begin{lstlisting}
  sudo blackman -a
\end{lstlisting}

\item Lista todas as categorias de ferramentas:
\begin{lstlisting}
  blackman -l
\end{lstlisting}

\item Lista uma categoria de ferramentas:
\begin{lstlisting}
  blackman -p categoria 
\end{lstlisting}

\end{itemize}

\subsection{Instalação pelo live-, netinstall- ISO ou ArchLinux}
Você pode instalar o BlackArch através de um live- ou netinstall-ISOs.\\Veja em
\url{https://www.blackarch.org/download.html#iso}.Os passos a seguir são executados depois da inicialização da ISO.
\begin{itemize}
\item Instalando o pacote blackarch-installer:
\begin{lstlisting}
  sudo pacman -S blackarch-installer
\end{lstlisting}

\item Execute
\begin{lstlisting}
  sudo blackarch-install
\end{lstlisting}

\end{itemize}

%------------------%
%  Chapter 3       %
%------------------%

\chapter{Guia do Desenvolvedor}

\section{Construir sistemas e repositórios Arch}

Os PKGBUILD são scripts de construição.Cada um diz ao makepkg(1) como construir o pacote.Os arquivos PKGBUILD são escritos em Bash.

Para mais informações, ler abaixo:
\begin{itemize}
\item \href{https://wiki.archlinux.org/index.php/Creating_Packages}{Arch Wiki: Creating Packages}
\item \href{https://wiki.archlinux.org/index.php/Makepkg}{Arch Wiki: makepkg}
\item \href{https://wiki.archlinux.org/index.php/PKGBUILD}{Arch Wiki: PKGBUILD}
\item \href{https://wiki.archlinux.org/index.php/Arch_Packaging_Standards}{Arch Wiki: Arch Packaging Standards}
\end{itemize}

\section{Padrão PKGBUILD do Blackarch}
Por sua simplicidade, os PKGBUILDs são similares aos do AUR, como algumas pequenas diferenças como mostrado a baixo.Cada pacote deve pertencer ao blackarch o minimo possível, e deve haver o máximo de crossover com vários pacotes percentes a vários grupos.
\subsection{Grupos}
Permite que os usuários instalem um específico tipo de pacote de forma rápida e fácil, os pacotes são separados por grupos.Os grupos permite que o usuário digite "pacman -S <group name>" e consiga vários pacotes.

\subsubsection{blackarch}
O grupo blackarch é a o grupo base onde todos os pacotes se encontram.Isto permite que o usuário instale todos os pacotes com facilidade.

O que deve estar aqui: Tudo.

\subsubsection{blackarch-anti-forensic}
Pacotes relacionados a o uso de contra forense, incluindo encriptação, estenografia, e tudo para modificar os atributos de arquivo/arquivos.
Inclui todas as ferramentas para trabalhar com tudo em geral para modificar um sistema com o propósito de esconder uma informação.

Exemplo: luks, TrueCrypt, Timestomp, dd, ropeadope, secure-delete

\subsubsection{blackarch-automation}
Pacotes para uso de ferramentas ou automatização de tarefas.

Exemplo: blueranger, tiger, wiffy

\subsubsection{blackarch-backdoor}
Pacotes para explorar ou abrir blackdoors em sistemas vulneráveis.

Exemplo: backdoor-factory, rrs, weevely

\subsubsection{blackarch-binary}
Pacotes para modificar arquivos binários de algum modo.

Exemplo: binwally, packerid

\subsubsection{blackarch-bluetooth}
Pacotes para explorar tudo relacionado ao Bluetooth padrão (802.15.1).

Exemplo: ubertooth, tbear, redfang

\subsubsection{blackarch-code-audit}
Pacotes para analisar um codigo em busca de vulnerabilidades.

Exemplo: flawfinder, pscan

\subsubsection{blackarch-cracker}
Pacotes usados para quebrar criptografia.

Exemplo: hashcat, john, crunch

\subsubsection{blackarch-crypto}
Pacotes para trabalhar com criptografia sem ser sua quebra.

Exemplo: ciphertest, xortool, sbd

\subsubsection{blackarch-database}
Pacotes que envolvam exploração de banco de dados em algum nível.

Exemplo: metacoretex, blindsql

\subsubsection{blackarch-debugger}
Pacotes para permitir que o usuário veja o que um programa em particular esta "fazendo" em tempo real.

Exemplo: radare2, shellnoob

\subsubsection{blackarch-decompiler}
Pacotes voltados para converter um programa compilado em codigo fonte.

Exemplo: flasm, jd-gui

\subsubsection{blackarch-defensive}
Pacotes para proteger o usuário de malware \& ataques de outros usuários.

Exemplos: arpon, chkrootkit, sniffjoke

\subsubsection{blackarch-disassembler}
Similar ao blackarch-decompiler, possui semelhança em alguns pacotes mas esses são mais voltados a produzir uma saída em assembly, não em codigo fonte.

Exemplo: inguma, radare2

\subsubsection{blackarch-dos}
Pacotes para o uso de ataques DoS (Denial of Service).

Exemplo: 42zip, nkiller2

\subsubsection{blackarch-drone}
Pacotes usados para o manejo físico de engenharia de drones.

Exemplo: meshdeck, skyjack

\subsubsection{blackarch-exploitation}
Pacotes para explorar outros programas ou serviços.

Exemplo: armitage, metasploit, zarp

\subsubsection{blackarch-fingerprint}
Pacotes para explorar equipamentos de biometria digital.

Exemplo: dns-map, p0f, httprint

\subsubsection{blackarch-firmware}
Pacotes para explorar vulnerabilidades no firmware

Exemplo: None yet, amend asap.

\subsubsection{blackarch-forensic}
Pacotes usados para procurar informação em discos físicos ou memoria.

Exemplo: aesfix, nfex, wyd

\subsubsection{blackarch-fuzzer}
Pacotes que usam o principio do teste fuzz, "jogando" números aleatórios com o objetivo de ver o que acontece.

Exemplo: msf, mdk3, wfuzz

\subsubsection{blackarch-hardware}
Pacotes que explorar ou manejam hardware.

Exemplo: arduino, smali

\subsubsection{blackarch-honeypot}
Pacotes que agem como "honeypots" (pote de mel), programas para simular uma vulnerabilidade e cria uma armadilha para os hackers.

Exemplo: artillery, bluepot, wifi-honey

\subsubsection{blackarch-keylogger}
Pacotes para gravar o digito em outros sistemas.

Exemplo: None yet, amend asap.

\subsubsection{blackarch-malware}
Pacotes para qualquer tipo de software malicioso ou detecção de
malware.

Exemplo: malwaredetect, peepdf, yara

\subsubsection{blackarch-misc}
Pacotes com nada em particular para encaixar em uma categoria.

Exemplo: oh-my-zsh-git, winexe, stompy

\subsubsection{blackarch-mobile}
Pacotes para manipulação de sistemas moveis.

Exemplo: android-sdk-platform-tools, android-udev-rules

\subsubsection{blackarch-networking}
Pacotes que involvem uma rede IP.

Exemplo: Praticamente tudo

\subsubsection{blackarch-nfc}
Pacotes que usam nfc (near-field communications).

Exemplo: nfcutils

\subsubsection{blackarch-packer}
Pacotes que operam ou envolvem packers.

/textif{packers são programas que juntam malware com outros executáveis.}

Exemplo: packerid

\subsubsection{blackarch-proxy}
Pacotes que agem com proxy, como redirecionar o tráfego para outro nó na internet.

Exemplo: burpsuite, ratproxy, sslnuke

\subsubsection{blackarch-recon}
Pacotes que procuram exploits.Mais no grupo umbrella para pacotes similares.

Exemplo: canri, dnsrecon, netmask

\subsubsection{blackarch-reversing}
Este é um grupo umbrella para todo decompilador, disassembler ou programas similares.

Exemplo: capstone, radare2, zerowine

\subsubsection{blackarch-scanner}
Pacotes que scaneiam um sistema em busca de vulnerabilidades.

Exemplo: scanssh, tiger, zmap

\subsubsection{blackarch-sniffer}
Pacotes que envolvem analise de tráfego de rede.

Exemplo: hexinject, pytactle, xspy

\subsubsection{blackarch-social}
Pacotes para ataque em redes sociais.

Exemplo: jigsaw, websploit

\subsubsection{blackarch-spoof}
Pacotes para falsificar um ataque, onde um ataque não é mostrado como uma ameaça a vitima.

Exemplo: arpoison, lans, netcommander

\subsubsection{blackarch-threat-model}
Pacotes que são usados para reportar/gravar um modelo de ameaça em um cenário particular.

Exemplo: magictree

\subsubsection{blackarch-tunnel}
Pacotes que são usados para criar tuneis no tráfego de rede.

Exemplo: ctunnel, iodine, ptunnel

\subsubsection{blackarch-unpacker}
Pacotes que são usados para extrair um malware de um executável.

Exemplo: js-beautify

\subsubsection{blackarch-voip}
Pacotes que operam em programas e protocolos voip.

Exemplo: iaxflood, rtp-flood, teardown

\subsubsection{blackarch-webapp}
Pacotes que operam em aplicações na internet.

Exemplo: metoscan, whatweb, zaproxy

\subsubsectiodo {blackarch-windows}
Este grupo é para todo pacote nativo para windows que é executado através do wine.

Exemplo: 3proxy-win32, pwdump, winexe

\subsubsection{blackarch-wireless}
Pacotes que operam em rede wireless em algum nível.

Exemplo: airpwn, mdk3, wiffy

\section{Estrutura do repositório}
Você pode achar o repositório git BlackArch principal aqui:
\href{https://github.com/BlackArch/blackarch}{https://github.com/BlackArch/blackarch}.
Temos também um secundário repositório aqui:
\href{https://github.com/BlackArch}{https://github.com/BlackArch}.

O repositório git principal, contém três diretórios importantes:

\begin{itemize}
\item docs - Documentação.
\item packages - Arquivo PKGBUILD.
\item scripts - Pequenos scripts.
\end{itemize}

\subsection{Scripts}
Aqui há uma referencia para o diretório de \verb|scripts/| scripts:

\begin{itemize}
\item baaur - Brevemente, irá subir os pacotes para o AUR.
\item babuild - Constrói os pacotes.
\item bachroot - Manipula o chroot para testes.
\item baclean - Limpa velhos arquivos .pkg.tar.xz do repositório.
\item baconflict - Irá substiuir o \verb|scripts/conflicts|.
\item bad-files - Acha arquivos ruins nos pacotes de construção.
\item balock - Obtém ou atualiza pacotes travados.
\item banotify - Notifica IRC sobre pushes de pacotes.
\item barelease - Atualiza pacotes no repositórios de pacotes.
\item baright - Imprime a informação de copyright do BlackArch.
\item basign - Assina pacotes.
\item basign-key - Assina as chaves.
\item blackman - Funciona de forma análoga ao pacman builds construindo pacotes do git(não confundir com o Blackman).
\item check-groups - Confere grupos.
\item checkpkgs - Confere erro em pacotes.
\item conflicts - Confere conflitos em pacotes.
\item dbmod - Modifica o banco de dados de pacotes.
\item depth-list - Cria uma lista de dependência.
\item deptree - Cria uma lista de dependência, somente dos pacotes do BlackArch.
\item get-blackarch-deps - Pega uma lista de dependências de pacotes do BlackArch.
\item get-official - Pega atualizações oficiais do BlackArch.
\item list-loose-packages - Lista os pacotes que estão ou não estão em um grupo que dependem de outros pacotes.
\item list-needed - Lista dependências perdidas.
\item list-removed - Lista os pacotes que estão no repositório mas não no git.
\item list-tools - Lista as ferramentas.
\item outdated - Olhar por pacotes desatualizados no repositório comparado ao git.
\item pkgmod - Modifica pacotes de construção.
\item pkgrel - Adiciona pkgrel em um pacote.
\item prep - Limpa arquivos PKGBUILD e procura erros.
\item sitesync - Sincroniza entre os pacotes locais e os pacotes no repositório.
\item size-hunt - Procura pacotes pesados/grandes.
\item source-backup - Realiza backup do codigo fonte dos arquivos.
\end{itemize}

\section{Contribuindo para o repositório}
Essa seção mostra como você pode contribuir com o projeto BlackArch Linux.
Nós aceitamos pull requests de todos os tamanhos, de pequenos consertos até novos pacotes.
\\Para mais ajuda, sugestões ou questões sinta-se livre para nos contatar.\\\\
Todo mundo que queira contribuir será bem vindo.Toda contribuição é bem vinda.

\subsection{Tutorial de requerimento}
Por favor leia este tutorial antes de contribuir:
\begin{itemize}
\item
\href{https://wiki.archlinux.org/index.php/Arch\_Packaging\_Standards)}{Arch
Packaging Standards}
\item \href{https://wiki.archlinux.org/index.php/Creating\_Packages}{Creating
Packages}
\item \href{https://wiki.archlinux.org/index.php/PKGBUILD}{PKGBUILD}
\item \href{https://wiki.archlinux.org/index.php/Makepkg}{Makepkg}
\end{itemize}

\subsection{Passos para contribuir}
Quando submeter suas mudanças para o projeto BlackArchLinux, siga estes passos:
\begin{enumerate}
\item Fork o repositório em
\url{https://github.com/BlackArch/blackarch}
\item Hackei os arquivos necessários, (e.g. PKGBUILD, .patch files, etc).
\item Commit as mudanças.
\item Push suas mudanças.
\item Nos peça para dar merge em suas mudanças, de preferencialmente um pull request.
\end{enumerate}

\subsection{Exemplo}
O exemplo a seguir demonstra como enviar um novo pacotes para o projeto BlackArch.
Nós usamos \href{https://wiki.archlinux.org/index.php/yaourt}{yaourt}
(you can use pacaur as well) para buscar arquivos PKGBUILD de
\textbf{nfsshell} para \href{https://aur.archlinux.org/}{AUR} e modificamos de acordo com a necessidade.

\subsubsection{Buscando PKGBUILD}
Busque o arquivo \textit{PKGBUILD} usando o yaourt ou pacaur:
\begin{lstlisting}
  user@blackarchlinux $ yaourt -G nfsshell
  ==> Download nfsshell sources
  x LICENSE
  x PKGBUILD
  x gcc.patch
  user@blackarchlinux $ cd nfsshell/
\end{lstlisting}

\subsubsection{Limpar o PKGBUILD}
Limpe o arquivo \textit{PKGBUILD} e salve:
\begin{lstlisting}
  user@blackarchlinux nfsshell $ ./blarckarch/scripts/prep PKGBUILD
  cleaning 'PKGBUILD'...
  expanding tabs...
  removing vim modeline...
  removing id comment...
  removing contributor and maintainer comments...
  squeezing extra blank lines...
  removing '|| return'...
  removing leading blank line...
  removing $pkgname...
  removing trailing whitespace...
\end{lstlisting}

\subsubsection{Modificações PKGBUILD}
Modifique o arquivo \textit{PKGBUILD}:
\begin{lstlisting}
  user@blackarchlinux nfsshell $ vi PKGBUILD
\end{lstlisting}

\subsubsection{Construa o pacote}
Construindo o pacote:
\begin{lstlisting}user@blackarchlinux nfsshell $ makepkg -sf
==> Making package: nfsshell 19980519-1 (Mon Dec  2 17:23:51 CET 2013)
==> Checking runtime dependencies...
==> Checking buildtime dependencies...
==> Retrieving sources...
-> Downloading nfsshell.tar.gz...
% Total    % Received % Xferd  Average Speed   Time    Time     Time
CurrentDload  Upload   Total   Spent    Left  Speed100 29213  100 29213    0
0  48150      0 --:--:-- --:--:-- --:--:-- 48206
-> Found gcc.patch
-> Found LICENSE
...
<lots of build process and compiler output here>
...
==> Leaving fakeroot environment.
==> Finished making: nfsshell 19980519-1 (Mon Dec  2 17:23:53 CET 2013)
\end{lstlisting}

\subsubsection{Instalando e testando os pacotes}
Intale e teste os pacotes:
\begin{lstlisting}
  user@blackarchlinux nfsshell $ pacman -U nfsshell-19980519-1-x86_64.pkg.tar.xz
  user@blackarchlinux nfsshell $ nfsshell # test it
\end{lstlisting}

\subsubsection{Adiciona, commit e push o pacote}
\begin{lstlisting}user@blackarchlinux nfsshell $ cd /blackarchlinux/packages
user@blackarchlinux ~/blackarchlinux/packages $ mv ~/nfsshell .
user@blackarchlinux ~/blackarchlinux/packages $ git commit -am nfsshell && git push
\end{lstlisting}

\subsubsection{Crie um pull request}
Create a pull request on \href{https://github.com/}{github.com}
\begin{lstlisting}
  firefox https://github.com/<contributor>/blackarchlinux
\end{lstlisting}

\subsubsection{Adicione remotamente}
Se melhor a se fazer se você trabalha remotamente é realizar o fork e o seu pull request para adicionar no repositório principal.
\begin{lstlisting}
  user@blackarchlinux ~/blackarchlinux $ git remote -v
  origin <the url of your fork> (fetch)
  origin <the url of your fork> (push)
  user@blackarchlinux ~/blackarchlinux $ git remote add upstream https://github.com/blackarch/blackarch
  user@blackarchlinux ~/blackarchlinux $ git remote -v
  origin <the url of your fork> (fetch)
  origin <the url of your fork> (push)
  upstream https://github.com/blackarch/blackarch (fetch)
  upstream https://github.com/blackarch/blackarch (push)
\end{lstlisting}

Por padrão o git deve realizar o push para a origem, mas tenha certeza que seu git
está configurado corretamente.Isto não deve causar problemas ao menos que você tenhas
alguns direitos de commit com não habilitar o push sem o mesmo.

Se você tem a habilidade de commitar, é mais fácil usar o
\textit{git@github.com:blackarch/blackarch.git} mas ai é com você.

\subsection{Requisitações}
\begin{enumerate}
\item Não adicione \textbf{Maintainer} ou \textbf{Contributor} comentarios nos arquivos
\textit{PKGBUILD}. Adicione o nome do mantedor ou contribuinte na seção de
AUTHORS do guia BlackArch.
\item Procure ser consistente, por favor siga o estilo dos outros arquivos \textit{PKGBUILD}
no repositório e use dois espaços para identação.
\end{enumerate}

\subsection{Dicas gerais}
\href{http://wiki.archlinux.org/index.php/Namcap}{namcap} para conferir os pacotes com erros.

%------------------%
%  Chapter 4       %
%------------------%

\chapter{Guia das ferramentas}
Em breve...

\section{Em breve}
Em breve...

%%% APPENDIX %%%
\appendix
%%%%%%%%%%%%%%%%%%%%%%%%%%%%%%%%%%%%%%%%%%%%%%%%%%%%%%%%%%%%%%%%%%%%%%%%%%%%%%%%
%                                                                              %
% BlackArch Linux appendix template                                            %
%                                                                              %
%%%%%%%%%%%%%%%%%%%%%%%%%%%%%%%%%%%%%%%%%%%%%%%%%%%%%%%%%%%%%%%%%%%%%%%%%%%%%%%%

\appendix

\chapter{Appendix}

\section{FAQs}

\section{AUTHORS}
The following people have contributed directly to BlackArch:
\begin{itemize}
\item Tyler Bennnett (tylerb@trix2voip.com)
\item fnord0 (fnord0@riseup.net)
\item nrz (nrz@nullsecurity.net)
\item Ellis Kenyo (elken.tdos@gmail.com)
\item CaledoniaProject (the.warl0ck.1989@gmail.com)
\item sudokode (sudokode@gmail.com)
\item Valentin Churavy (v.churavy@gmail.com)
\item Boy Sandy Gladies Arriezona (reno.esper@gmail.com)
\item Mathias Nyman
\item Johannes Löthberg (demizide@gmail.com)
\item Thiago da Silva Teixeira (teixeira.zeus@gmail.com)
\end{itemize}
The following people contributed directly to ArchPwn, which has been merged
into BlackArch:
\begin{itemize}
\item Francesco Piccinno (stack.box@gmail.com)
\item jensp (jens@jenux.homelinux.org)
\item Valentin Churavy (v.churavy@gmail.com)
\end{itemize}
We have taken build code from the following people:
\begin{itemize}
\item  3ED (krzysztof1987@gmail.com)
\item  AUR Perl (aurperl@juster.info)
\item  Aaron Griffin (aaron@archlinux.org)
\item  Abakus (java5@arcor.de)
\item  Adam Wolk (netprobe@gmail.com)
\item  Aleix Pol (aleixpol@kde.org)
\item  Aleshus (aleshusi@gmail.com)
\item  Alessandro Pazzaglia (jackdroido@gmail.com)
\item  Alessandro Sagratini (ale\_sagra@hotmail.com)
\item  Alex Cartwright (alexc223@googlemail.com)
\item  Alexander De Sousa (archaur.xandy21@spamgourmet.com)
\item  Alexander Rødseth (rodseth@gmail.com)
\item  Allan McRae (allan@archlinux.org)
\item  AmaN (gabroo.punjab.da@gmail.com)
\item  Andre Klitzing (aklitzing@online.de)
\item  Andrea Scarpino (andrea@archlinux.org)
\item  Andreas Schönfelder (passtschu@freenet.de)
\item  Andrej Gelenberg (andrej.gelenberg@udo.edu)
\item  Angel Velasquez (angvp@archlinux.org)
\item  Antoine Lubineau (antoine@lubignon.info)
\item  Anton Bazhenov (anton.bazhenov@gmail.com)
\item  Arkham (arkham@archlinux.us)
\item  Arthur Danskin (arthurdanskin@gmail.com)
\item  Balda (balda@balda.ch)
\item  Balló György (ballogyor+arch@gmail.com)
\item  Bartek Piotrowski (barthalion@gmail.com)
\item  Bartosz Feński (fenio@debian.org)
\item  Bartłomiej Piotrowski (nospam@bpiotrowski.pl)
\item  Bogdan Szczurek (thebodzio@gmail.com)
\item  Brad Fanella (bradfanella@archlinux.us)
\item  Brian Bidulock (bidulock@openss7.org)
\item  C Anthony Risinger (anthony@xtfx.me)
\item  CRT (crt.011@gmail.com)
\item  Can Celasun (dcelasun@gmail.com)
\item  Chaniyth (chaniyth@yahoo.com)
\item  Chris Brannon (cmbrannon79@gmail.com)
\item  Chris Giles (Chris.G.27@gmail.com) \& daschu117
\item  Christoph Siegenthaler (csi@gmx.ch)
\item  Christoph Zeiler (archNOSPAM@moonblade.org)
\item  Clément DEMOULINS (clement@archivel.fr)
\item  Corrado Primier (bardo@aur.archlinux.org)
\item  Daenyth (Daenyth+Arch@gmail.com)
\item  Dale Blount (dale@archlinux.org)
\item  Damir Perisa (damir.perisa@bluewin.ch)
\item  Dan Fuhry (dan@fuhry.us)
\item  Dan Serban (dserban01@yahoo.com)
\item  Daniel A. Campoverde Carrión
\item  Daniel Golle
\item  Daniel Griffiths (ghost1227@archlinux.us)
\item  Daniel J Griffiths (ghost1227@archlinux.us)
\item  Daniel Micay (danielmicay@gmail.com)
\item  Dave Reisner (dreisner@archlinux.org)
\item  Dawid Wrobel (cromo@klej.net)
\item  Devaev Maxim (mdevaev@gmail.com)
\item  Devin Cofer (ranguvar@archlinux.us)
\item  DigitalPathogen (aur@InfoSecResearchLabs.co.uk)
\item  DigitalPathogen (aur@digitalpathogen.co.uk)
\item  Dmitry A. Ilyashevich (dmitry.ilyashevich@gmail.com)
\item  Dominik Heidler (dheidler@gmail.com)
\item  DrZaius (lou@fakeoutdoorsman.com)
\item  Ebubekir KARUL (ebubekirkarul@yandex.com)
\item  Eduard "bekks" Warkentin (eduard.warkentin@gmail.com)
\item  Elmo Todurov (todurov@gmail.com)
\item  Emmanuel Gil Peyrot (linkmauve@linkmauve.fr)
\item  Eric Belanger (eric@archlinux.org)
\item  Ermak (ermak@email.it)
\item  Evangelos Foutras (evangelos@foutrelis.com)
\item  Fabian Melters (melters@gmail.com)
\item  Fabiano Furtado (fusca14@gmail.com)
\item  Federico Quagliata (ntp@quaqo.org)
\item  Firmicus (francois.archlinux@org)
\item  Florian Pritz (bluewind@jabber.ccc.de)
\item  Florian Pritz (flo@xinu.at)
\item  Francesco Piccinno (stack.box@gmail.com)
\item  François Charette (francois@archlinux.org)
\item  Gaetan Bisson (bisson@archlinux.org)
\item  Geoffroy Carrier (geoffroy.carrier@koon.fr)
\item  Georg Grabler (STiAT)
\item  George Hilliard (gh403@msstate.edu)
\item  Gerardo Exequiel Pozzi (vmlinuz386@yahoo.com.ar)
\item  Gilles CHAUVIN (gcnweb@gmail.com)
\item  Giovanni Scafora (giovanni@archlinux.org)
\item  Gordin (9ordin@gmail.com)
\item  Guillaume ALAUX (guillaume@archlinux.org)
\item  Guillermo Ramos (0xwille@gmail.com)
\item  Gustavo Alvarez (sl1pkn07@gmail.com)
\item  Hugo Doria (hugo@archlinux.org)
\item  Hyacinthe Cartiaux (hyacinthe.cartiaux@free.fr)
\item  James Fryman (jfryman@gmail.com)
\item  Jan "heftig" Steffens (jan.steffens@gmail.com)
\item  Jan de Groot (jgc@archlinux.org)
\item  Jaroslav Lichtblau (dragonlord@aur.archlinux.org)
\item  Jaroslaw Swierczynski (swiergot@aur.archlinux.org)
\item  Jason Chu (jason@archlinux.org)
\item  Jason R Begley (jayray@digitalgoat.com)
\item  Jason Rodriguez
\item  Jason St. John (jstjohn@purdue.edu)
\item  Jawmare (victor2008@gmail.com)
\item  Jeff Mickey (jeff@archlinux.org)
\item  Jens Pranaitis (jens@chaox.net)
\item  Jens Pranaitis (jens@jenux.homelinux.org)
\item  Jinx (jinxware@gmail.com)
\item  John D Jones III (jnbek1972@gmail.com)
\item  John Proctor (jproctor@prium.net)
\item  Jon Bergli Heier (snakebite@jvnv.net)
\item  Jonas Heinrich
\item  Jonathan Steel (jsteel@aur.archlinux.org)
\item  Joris Steyn (jorissteyn@gmail.com)
\item  Josh VanderLinden (arch@cloudlery.com)
\item  Jozef Riha (jose1711@gmail.com)
\item  Judd Vinet (jvinet@zeroflux.org)
\item  Juergen Hoetzel (jason@archlinux.org)
\item  Juergen Hoetzel (juergen@archlinux.org)
\item  Justin Davis (jrcd83@gmail.com)
\item  Kaiting Chen (kaitocracy@gmail.com)
\item  Kaos
\item  Kevin Piche (kevin@archlinux.org)
\item  Kory Woods (kory@virlo.net)
\item  Kyle Keen (keenerd@gmail.com)
\item  Larry Hajali (larryhaja@gmail.com)
\item  LeCrayonVert
\item  Le\_suisse (lesuisse.dev+aur@gmail.com)
\item  Lekensteyn (lekensteyn@gmail.com)
\item  Limao Luo (luolimao+AUR@gmail.com)
\item  Lucien Immink
\item  Lukas Fleischer (archlinux@cryptocrack.de)
\item  Manolis Tzanidakis
\item  Marcin "avalan" Falkiewicz (avalatron@gmail.com)
\item  Mariano Verdu (verdumariano@gmail.com)
\item  Marti Raudsepp (marti@juffo.org)
\item  MatToufoutu (mattoufootu@gmail.com)
\item  Matthew Sharpe (matt.sharpe@gmail.com)
\item  Mauro Andreolini (mauro.andreolini@unimore.it)
\item  Max Pray a.k.a. Synthead (synthead@gmail.com)
\item  Max Roder (maxroder@web.de)
\item  Maxwell Pray a.k.a. Synthead (synthead@gmail.com)
\item  Maxwell Pray a.k.a. Synthead (synthead1@gmail.com)
\item  Mech (tiago.bmp@gmail.com)
\item  Michael Düll (mail@akurei.me)
\item  Michael P (ptchinster@archlinux.us)
\item  Michal Krenek (mikos@sg1.cz)
\item  Michal Zalewski (lcamtuf@coredump.cx)
\item  Miguel Paolino (mpaolino@gmail.com)
\item  Miguel Revilla (yo@miguelrevilla.com)
\item  Mike Roberts (noodlesgc@gmail.com)
\item  Mike Sampson (mike@sambodata.com)
\item  Nassim Kacha (nassim.kacha@gmail.com)
\item  Nicolas Pouillard (nicolas.pouillard@gmail.com)
\item  Nicolas Pouillard https://nicolaspouillard.fr
\item  Niklas Schmuecker
\item  Oleander Reis (oleander@oleander.cc)
\item  Olivier Le Moal (mail@olivierlemoal.fr)
\item  Olivier Médoc "oliv" (o\_medoc@yahoo.fr)
\item  Pascal E. (archlinux@hardfalcon.net)
\item  Patrick Leslie Polzer (leslie.polzer@gmx.net)
\item  Paul Mattal (paul@archlinux.org)
\item  Paul Mattal (pjmattal@elys.com)
\item  Pengyu CHEN (cpy.prefers.you@gmail.com)
\item  Peter Wu (lekensteyn@gmail.com)
\item  Philipp 'TamCore' B. (philipp@tamcore.eu)
\item  Pierre Schmitz (pierre@archlinux.de)
\item  Pranay Kanwar (pranay.kanwar@gmail.com)
\item  Pranay Kanwar (warl0ck@metaeye.org)
\item  PyroPeter (googlemail@com.abi1789)
\item  PyroPeter (googlemail.com@abi1789)
\item  Ray Rashif (schiv@archlinux.org)
\item  Remi Gacogne
\item  Renan Fernandes (renan@kauamanga.com)
\item  Richard Murri (admin@richardmurri.com)
\item  Roberto Alsina (ralsina@kde.org)
\item  Robson Peixoto (robsonpeixoto@gmail.com)
\item  Roel Blaauwgeers (roel@ttys0.nl)
\item  Rorschach (r0rschach@lavabit.com)
\item  Ruben Schuller (shiml@orgizm.net)
\item  Rudy Matela (rudy@matela.com)
\item  Ryon Sherman (ryon.sherman@gmail.com)
\item  Sabart Otto \item  Seberm (seberm@gmail.com)
\item  SakalisC (chrissakalis@gmail.com)
\item  Sam Stuewe (halosghost@archlinux.info)
\item  SanskritFritz (SanskritFritz@gmail.com)
\item  Sarah Hay (sarahhay@mb.sympatico)
\item  Sebastian Benvenuti (sebastianbenvenuti@gmail.com)
\item  Sebastian Nowicki (sebnow@gmail.com)
\item  Sebastien Duquette (ekse.0x@gmail.com)
\item  Sebastien LEDUC (sebastien@sleduc.fr)
\item  Sebastien Leduc (sebastien@sleduc.fr)
\item  Sergej Pupykin (pupykin.s+arch@gmail.com)
\item  Sergio Rubio (rubiojr@biondofu.net)
\item  Sheng Yu (magicfish1990@gmail.com)
\item  Simon Busch (morphis@gravedo.de)
\item  Simon Legner (Simon.Legner@gmail.com)
\item  Sirat18 (aur@sirat18.de)
\item  SpepS (dreamspepser@yahoo.it)
\item  Spider.007 (archPackage@spider007.net)
\item  Stefan Seering
\item  Stephane Travostino (stephane.travostino@gmail.com)
\item  Stéphane Gaudreault (stephane@archlinux.org)
\item  Sven Kauber (celeon@gmail.com)
\item  Sven Schulz (omee@archlinux.de)
\item  Sébastien Duquette (ekse.0x@gmail.com)
\item  Sébastien Luttringer (seblu@archlinux.org)
\item  TDY (tdy@gmx.com)
\item  Teemu Rytilahti (tpr@iki.fi)
\item  Testuser\_01
\item  Thanx (thanxm@gmail.com)
\item  Thayer Williams (thayer@archlinux.org)
\item  Thomas S Hatch (thatch45@gmail.com)
\item  Thorsten Töpper
\item  Tilmann Becker (tilmann.becker@freenet.de)
\item  Timothy Redaelli (timothy.redaelli@gmail.com)
\item  Timothée Ravier (tim@siosm.fr)
\item  Tino Reichardt
\item  Tobias Kieslich (tobias@justdreams.de)
\item  Tobias Powalowski (tpowa@archlinux.org)
\item  Tom K (tomk@runbox.com)
\item  Tom Newsom (Jeepster@gmx.co.uk)
\item  Tomas Lindquist Olsen (tomas@famolsen.dk)
\item  Travis Willard (travisw@wmpub.ca)
\item  Valentin Churavy (v.churavy@gmail.com)
\item  ViNS (gladiator@fastwebnet.it)
\item  Vlatko Kosturjak (kost@linux.hr)
\item  Wes Brown (wesbrown18@gmail.com)
\item  William Rea (sillywilly@gmail.com)
\item  Xavier Devlamynck (magicrhesus@ouranos.be)
\item  Xiao\item Long Chen (chenxiaolong@cxl.epac.to)
\item  aeolist (aeolist@hotmail.com)
\item  ality@pbrane.org
\item  astaroth (astaroth\_@web.de)
\item  bender02@archlinux.us
\item  billycongo (billycongo@gmail.com)
\item  bslackr (brendan@vastactive.com)
\item  cbreaker (cbreaker@tlen.pl)
\item  chimeracoder (dev@chimeracoder.net)
\item  damir (damir@archlinux.org)
\item  danitool
\item  darkapex (me@jailuthra.in)
\item  daronin
\item  dkaylor (dpkaylor@gmail.com)
\item  dobo (dobo90\_at\_gmail@com)
\item  dorphell (dorphell@archlinux.org)
\item  evr (evanroman@at.gmail)
\item  fnord0 (fnord0@riseup.net)
\item  fxbru (frxbru@gmail)
\item  hcar
\item  icarus (icarus.roaming@gmail.com)
\item  iceman (icemanf@gmail.com)
\item  kastor (kastor@fobos.org)
\item  kfgz (kfgz@interia.pl)
\item  linuxSEAT (linuxSEAT@gmail.com)
\item  m4xm4n (max@maxfierke.com)
\item  mar77i (mysatyre@gmail.com)
\item  marc0s (marc0s@fsfe.org)
\item  mickael9 (mickael9@gmail.com)
\item  nblock (nblock@archlinux.us)
\item  nofxx (x@nofxx.com)
\item  onny (onny@project
\item  pootzko (pootzko@gmail.com)
\item  revel (revel@muub.net)
\item  rich\_o (rich\_o@lavabit.com)
\item  s1gma (s1gma@mindslicer.com)
\item  sandman (r.coded@gmail.com)
\item  sebikul (sebikul@gmail.com)
\item  sh0 (mee@sh0.org)
\item  shild (sxp@bk.ru)
\item  simo (simo@archlinux.org)
\item  snuo
\item  sudokode (sudokode@gmail.com)
\item  tobias (tobias@archlinux.org)
\item  trashstar (trash@ps3zone.org)
\item  unexist (unexist@subforge.org)
\item  untitled (rnd0x00@gmail.com)
\item  virtuemood (virtue@archlinux.us)
\item  wido (widomaker2k7@gmail.com)
\item  wodim (neikokz@gmail.com)
\item  yannsen (ynnsen@gmail.com)
\end{itemize}


\end{document}

%%% EOF %%%
